% =====================================================================
\chapter{Soluciones de los ejercicios}
\label{ap:soluciones}
% =====================================================================

Cada solución se presenta en cinco partes: datos, planteamiento, desarrollo, resultado y comprobación. En los pasos intermedios se redondea a cuatro cifras significativas salvo que se indique otra cosa. Cuando un logaritmo se multiplica después por un exponente grande o se eleva a una potencia, se conservan cinco decimales, porque un error en la cuarta cifra de $\log_{10}N$ se nota en la tercera cifra de $N$. Todas las cifras se han comprobado con un programa en Python.

% ---------------------------------------------------------------------
\section{Tema 1}

\begin{solucion}{eje:1-sesgo}
\datos clasificador entrenado con páginas citadas desde la Wikipedia en inglés como positivos; rastreo multilingüe.
\planteamiento el clasificador aprende a reconocer lo que se parece a sus positivos; todo lo que difiere de ellos por razones ajenas a la calidad recibe puntuación baja.
\pasos
\begin{enumerate}[label=\arabic*.]
  \item Texto en otros idiomas, sobre todo los que tienen pocas páginas citadas en la Wikipedia inglesa: el clasificador confunde «distinto del inglés enciclopédico» con «de baja calidad».
  \item Registros informales y variedades dialectales (foros, conversación, literatura popular), útiles para un modelo de lenguaje general y ausentes del estilo enciclopédico.
  \item Consecuencia: el modelo resultante ve menos de esas variedades y las modeliza peor; su perplejidad sobre ellas es mayor y sus respuestas se sesgan hacia un registro y una cultura concretos.
  \item Medidas: aplicar el clasificador por idioma, con positivos de referencia de cada idioma; fijar el umbral por idioma o por dominio; y auditar muestras de lo descartado.
\end{enumerate}
\resultado \textbf{se descartan injustificadamente el texto en idiomas con pocos recursos y los registros informales o dialectales; el modelo hereda ese sesgo. Se mitiga con clasificadores y umbrales por idioma y con auditoría de lo descartado.}
\comprobacion la observación del \cref{tema:1} sobre la calidad como elección describe exactamente este mecanismo.
\end{solucion}

\begin{solucion}{eje:1-gopher}
\datos la tabla de cuatro documentos y las cuatro reglas de la \cref{def:gopher}.
\planteamiento calcular por documento la longitud media de palabra, la fracción de líneas con «...» y la fracción de palabras con letra; basta que falle una regla para descartarlo.
\pasos
\begin{enumerate}[label=\arabic*.]
  \item A: $40<50$ palabras, \textbf{falla} la regla~(i). El resto pasaría: $200/40=5{,}0$; $0/4=0$; $38/40=0{,}95$.
  \item B: $1200$ palabras, bien. Media $5640/1200=4{,}7$, bien. Líneas con «...»: $24/60=0{,}40>0{,}30$, \textbf{falla}. Letras: $1170/1200=0{,}975$.
  \item C: $800$ palabras; media $4160/800=5{,}2$; «...» $6/50=0{,}12\le 0{,}30$; letras $720/800=0{,}90\ge 0{,}80$. \textbf{Pasa}.
  \item D: $3000$ palabras; media $13\,500/3000=4{,}5$; «...» $5/100=0{,}05$; letras $2250/3000=0{,}75<0{,}80$, \textbf{falla}.
\end{enumerate}
\resultado \textbf{solo pasa C.} A es demasiado corto, B parece un listado de fragmentos truncados y D tiene demasiadas «palabras» sin letras (tablas, precios, códigos).
\comprobacion cada documento descartado falla exactamente una regla y con margen claro ($40$ frente a $50$; $0{,}40$ frente a $0{,}30$; $0{,}75$ frente a $0{,}80$): ningún resultado depende del redondeo.
\end{solucion}

\begin{solucion}{eje:1-mezcla}
\datos $D=10^{12}$. Web: $U=8\cdot 10^{11}$, $w=0{,}60$. Código: $2\cdot 10^{11}$, $0{,}25$. Libros: $5\cdot 10^{10}$, $0{,}10$. Ciencia: $10^{10}$, $0{,}05$.
\planteamiento $e_i=w_iD/U_i$ (\cref{def:mezcla}); más de unas $4$ épocas rinde poco (\cref{tema:5}).
\pasos
\begin{enumerate}[label=\arabic*.]
  \item Web: $0{,}60\cdot 10^{12}=6\cdot 10^{11}$ tokens vistos; $6\cdot 10^{11}/8\cdot 10^{11}=0{,}75$ épocas.
  \item Código: $2{,}5\cdot 10^{11}/2\cdot 10^{11}=1{,}25$ épocas.
  \item Libros: $10^{11}/5\cdot 10^{10}=2$ épocas.
  \item Ciencia: $5\cdot 10^{10}/10^{10}=5$ épocas.
  \item Ajuste: con $w=0{,}04$, la ciencia queda en $4\cdot 10^{10}/10^{10}=4$ épocas; el $0{,}01$ sobrante pasa a la web, que sube a $0{,}61\cdot 10^{12}/8\cdot 10^{11}=0{,}7625$ épocas.
\end{enumerate}
\resultado \textbf{web $0{,}75$; código $1{,}25$; libros $2$; ciencia $5$ épocas. Preocupa la ciencia ($5>4$); bajar su peso a $0{,}04$ y pasar $0{,}01$ a la web la deja en $4$.}
\comprobacion los pesos suman $0{,}60+0{,}25+0{,}10+0{,}05=1$ y los tokens vistos, $6+2{,}5+1+0{,}5=10$ (en $10^{11}$), es decir, $10^{12}=D$.
\end{solucion}

\begin{solucion}{eje:1-ppl}
\datos $q(\text{de})=0{,}4$, $q(\text{la})=0{,}3$, $q(\text{casa})=0{,}2$, $q(\text{oferta})=0{,}1$; documentos A, B y C de cuatro palabras; umbral $\PPL\le 4$.
\planteamiento $\PPL_q(x)=\bigl(\prod_tq(x_t)\bigr)^{-1/T}$ con $T=4$; para A, comparar $\ln\PPL$ con $H(\hat p)+D(\hat p\,\|\,q)$.
\pasos
\begin{enumerate}[label=\arabic*.]
  \item A: $\prod q=0{,}3\cdot 0{,}2\cdot 0{,}4\cdot 0{,}3=0{,}0072$; $\ln 0{,}0072=-4{,}9337$; $\ln\PPL=4{,}9337/4=1{,}2334$; $\PPL=e^{1{,}2334}=3{,}433$.
  \item B: $\prod q=0{,}1^{3}\cdot 0{,}2=0{,}0002$; $\ln\PPL=8{,}5172/4=2{,}1293$; $\PPL=8{,}409$.
  \item C: $\prod q=0{,}4^{4}$; $\PPL=1/0{,}4=2{,}5$.
  \item Con umbral~$4$ se conservan A y C, y se descarta B.
  \item Comprobación de la \cref{prop:ppl-kl} para A: $\hat p=(\text{de}:1/4,\ \text{la}:1/2,\ \text{casa}:1/4)$. $H(\hat p)=-\bigl(\tfrac14\ln\tfrac14+\tfrac12\ln\tfrac12+\tfrac14\ln\tfrac14\bigr)=1{,}0397$. $D(\hat p\,\|\,q)=\tfrac14\ln\tfrac{0{,}25}{0{,}4}+\tfrac12\ln\tfrac{0{,}5}{0{,}3}+\tfrac14\ln\tfrac{0{,}25}{0{,}2}=-0{,}1175+0{,}2554+0{,}0558=0{,}1937$. Suma: $1{,}2334$.
\end{enumerate}
\resultado \textbf{$\PPL(\text{A})=3{,}433$, $\PPL(\text{B})=8{,}409$, $\PPL(\text{C})=2{,}5$; pasan A y C.} El documento degenerado C obtiene la mejor puntuación, como predice el \cref{cor:repeticion}.
\comprobacion $1/q_{\max}=1/0{,}4=2{,}5$ es la cota inferior del \cref{cor:repeticion}, y la alcanza exactamente C; $H(\hat p)+D(\hat p\,\|\,q)=1{,}2334=\ln 3{,}433$.
\end{solucion}

\begin{solucion}{eje:1-fermi}
\datos $2500$ archivos en $1{,}5$ h por máquina; $100\,000$ archivos por volcado; $96$ volcados.
\planteamiento el filtrado es paralelo por archivos: el tiempo es proporcional al número de archivos e inversamente proporcional al número de máquinas.
\pasos
\begin{enumerate}[label=\arabic*.]
  \item (a) $100\,000/2500=40$ lotes; $40\cdot 1{,}5=60$ horas en una máquina.
  \item (b) $60/2=30$ máquinas.
  \item (c) $96\cdot 60=5760$ horas-máquina.
\end{enumerate}
\resultado \textbf{(a) $60$ h; (b) $30$ máquinas; (c) $5760$ horas-máquina.}
\comprobacion $30$ máquinas $\times\,2$ h $=60$ horas-máquina, el coste de un volcado; el orden de magnitud explica por qué la práctica de CS336 insiste en paralelizar~\cite[§4]{cs336a4}.
\end{solucion}

\begin{solucion}{eje:1-minmax}
\datos $D$ fijo; fuentes con $U_i$ tokens únicos; $e_i=w_iD/U_i$; tabla del \cref{ej:llama-mezcla}.
\planteamiento acotar inferiormente $\max_ie_i$ sumando las desigualdades $e_i\le\max_je_j$, y estudiar cuándo se alcanza la cota.
\pasos
\begin{enumerate}[label=\arabic*.]
  \item Sea $e=\max_ie_i$. Para todo $i$, $w_iD/U_i\le e$, es decir, $w_i\le eU_i/D$.
  \item Sumando y usando $\sum_iw_i=1$: $1\le e\sum_iU_i/D$, luego $e\ge D/\sum_iU_i$ para toda mezcla.
  \item La mezcla proporcional $w_i=U_i/\sum_jU_j$ da $e_i=D/\sum_jU_j$ para todo $i$: alcanza la cota y es óptima.
  \item Unicidad: si una mezcla alcanza $e=D/\sum_jU_j$, la suma del paso~2 es una igualdad, $\sum_iw_i=\sum_ieU_i/D$, y como cada sumando cumple $w_i\le eU_i/D$, todas las desigualdades son igualdades: $w_i=U_i/\sum_jU_j$.
  \item LLaMA: $\sum_iU_i=852{,}7+198{,}1+98{,}4+25{,}7+28{,}3+33{,}0+27{,}2=1263{,}4$ (en $10^{9}$). El mínimo es $1400/1263{,}4=1{,}108$ épocas.
\end{enumerate}
\resultado \textbf{la mezcla proporcional es la única que minimiza $\max_ie_i$, con valor $D/\sum_iU_i$; para LLaMA, $1{,}108$ épocas, frente a las $2{,}45$ que sus autores eligieron para Wikipedia.}
\comprobacion la elección de LLaMA no minimiza la repetición: sobremuestrea a propósito las fuentes pequeñas y valiosas, aceptando hasta $2{,}45$ épocas, por debajo del límite de unas $4$ del \cref{tema:5}.
\end{solucion}

% ---------------------------------------------------------------------
\section{Tema 2}

\begin{solucion}{eje:2-porque}
\datos dos corpus del mismo tamaño en tokens, uno con muchos duplicados y otro deduplicado.
\planteamiento comparar cuánta información nueva aporta cada token y qué hace el modelo con el texto que ve muchas veces.
\pasos
\begin{enumerate}[label=\arabic*.]
  \item \emph{Cómputo malgastado.} Un token duplicado no aporta información nueva: la entropía cruzada sobre él ya se redujo la primera vez. Con $C\approx 6ND$ fijo, el corpus efectivo es menor que su tamaño nominal.
  \item \emph{Memorización.} Lo repetido cientos de veces se aprende literalmente y puede regurgitarse, incluidos datos personales o protegidos~\cite{lee}.
  \item Además, \emph{contaminación}: si los duplicados cruzan la frontera entre entrenamiento y evaluación, las métricas se inflan.
\end{enumerate}
\resultado \textbf{(i) desperdicio de cómputo en información repetida; (ii) memorización literal (y, además, contaminación de la evaluación).}
\comprobacion caso límite: un corpus formado por un único documento repetido un millón de veces tiene el mismo tamaño en tokens, pero solo permite aprender ese documento.
\end{solucion}

\begin{solucion}{eje:2-conjuntos}
\datos $S_1=\{a,b,c,d\}$, $S_2=\{b,c,e\}$, $S_3=\{a,c,e,f\}$.
\planteamiento $J=\abs{X\cap Y}/\abs{X\cup Y}$ con $\abs{X\cup Y}=\abs X+\abs Y-\abs{X\cap Y}$.
\pasos
\begin{enumerate}[label=\arabic*.]
  \item $S_1,S_2$: intersección $\{b,c\}$; unión $4+3-2=5$; $J=2/5=0{,}4$.
  \item $S_1,S_3$: intersección $\{a,c\}$; unión $4+4-2=6$; $J=2/6=1/3$.
  \item $S_2,S_3$: intersección $\{c,e\}$; unión $3+4-2=5$; $J=2/5=0{,}4$.
\end{enumerate}
\resultado \textbf{$J(S_1,S_2)=0{,}4$; $J(S_1,S_3)=1/3\approx 0{,}333$; $J(S_2,S_3)=0{,}4$. Empatan $S_1$--$S_2$ y $S_2$--$S_3$.}
\comprobacion las tres intersecciones tienen $2$ elementos; el par con mayor unión, $S_1$--$S_3$, debe tener la menor similitud, y así es.
\end{solucion}

\begin{solucion}{eje:2-luna}
\datos A~=~«la luna brilla sobre el mar», B~=~«la luna brilla sobre el río»; $k=2$.
\planteamiento un texto de $T$ palabras tiene $T-1$ shingles de $2$ palabras; $J=\abs{A\cap B}/\abs{A\cup B}$.
\pasos
\begin{enumerate}[label=\arabic*.]
  \item A: $\{$la luna, luna brilla, brilla sobre, sobre el, el mar$\}$, $5$ shingles.
  \item B: $\{$la luna, luna brilla, brilla sobre, sobre el, el río$\}$, $5$ shingles.
  \item Intersección: los cuatro primeros, $\abs{A\cap B}=4$. Unión: $5+5-4=6$.
  \item $J=4/6=2/3$.
\end{enumerate}
\resultado \textbf{$J=2/3\approx 0{,}667$.}
\comprobacion solo cambia la última palabra, que aparece en un único shingle de cada frase: cada conjunto pierde $1$ de sus $5$ elementos y gana otro distinto. La fracción compartida de cada conjunto es $4/5$, pero la unión tiene $6>5$ elementos, de modo que $J<4/5$.
\end{solucion}

\begin{solucion}{eje:2-firmas}
\datos $S_1=\{0,2,5\}$, $S_2=\{1,2,3\}$, $S_3=\{0,3,5\}$; $h_1$, $h_2$, $h_3$ del enunciado.
\planteamiento una función es permutación del universo si toma los $6$ valores distintos; la firma es $(\min h_1,\min h_2,\min h_3)$ sobre cada conjunto; la similitud estimada es la fracción de coordenadas iguales.
\pasos
\begin{enumerate}[label=\arabic*.]
  \item Valores en $x=0,\dots,5$: $h_1=(1,2,3,4,5,0)$, permutación; $h_2=(1,0,5,4,3,2)$, permutación ($5$ y $6$ son primos entre sí); $h_3=(3,5,1,3,5,1)$, no lo es ($2$ y $6$ comparten el factor~$2$).
  \item $S_1=\{0,2,5\}$: $\min(1,3,0)=0$; $\min(1,5,2)=1$; $\min(3,1,1)=1$. Firma $(0,1,1)$.
  \item $S_2=\{1,2,3\}$: $\min(2,3,4)=2$; $\min(0,5,4)=0$; $\min(5,1,3)=1$. Firma $(2,0,1)$.
  \item $S_3=\{0,3,5\}$: $\min(1,4,0)=0$; $\min(1,4,2)=1$; $\min(3,3,1)=1$. Firma $(0,1,1)$.
  \item Similitudes reales: $J(S_1,S_2)=1/5=0{,}2$; $J(S_1,S_3)=2/4=0{,}5$; $J(S_2,S_3)=1/5=0{,}2$.
  \item Estimadas con las tres funciones: $S_1$--$S_2$ coinciden solo en $h_3$, $1/3$; $S_1$--$S_3$ en las tres, $1$; $S_2$--$S_3$ solo en $h_3$, $1/3$. Solo con $h_1$ y $h_2$: $0$, $1$ y $0$.
\end{enumerate}
\resultado \textbf{(a) $h_1$ y $h_2$ son permutaciones; $h_3$ no. (b) $\sig(S_1)=(0,1,1)$, $\sig(S_2)=(2,0,1)$, $\sig(S_3)=(0,1,1)$. (c) Reales: $0{,}2$; $0{,}5$; $0{,}2$. Estimadas: $1/3$; $1$; $1/3$.} Con tan pocas funciones el error es grande, y $h_3$, que no es permutación, coincide más de la cuenta.
\comprobacion por el \cref{cor:estimador}, con $J=0{,}5$ y $n=3$ la desviación típica es $\sqrt{0{,}25/3}=0{,}29$: estimar~$1$ está a menos de dos desviaciones típicas y es compatible con el azar.
\end{solucion}

\begin{solucion}{eje:2-lsh}
\datos $r=4$, $b=25$; $s=0{,}3$ y $s=0{,}85$.
\planteamiento $P(s)=1-(1-s^{r})^{b}$ (\cref{prop:curva-s}); para elevar a $25$, $x^{25}=e^{25\ln x}$.
\pasos
\begin{enumerate}[label=\arabic*.]
  \item $s=0{,}3$: $s^{4}=0{,}0081$; $\ln(0{,}9919)=-0{,}008133$; $25\cdot(-0{,}008133)=-0{,}2033$; $e^{-0{,}2033}=0{,}8160$; $P=0{,}1840$.
  \item $s=0{,}85$: $0{,}85^{4}=0{,}5220$; $\ln 0{,}4780=-0{,}7382$; $25\cdot(-0{,}7382)=-18{,}45$; $e^{-18{,}45}=9{,}7\cdot 10^{-9}$; $P=1-9{,}7\cdot 10^{-9}$.
  \item Umbral aproximado: $(1/25)^{1/4}=\sqrt{\sqrt{0{,}04}}=\sqrt{0{,}2}=0{,}4472$.
  \item Inflexión exacta: $\bigl((r-1)/(rb-1)\bigr)^{1/r}=(3/99)^{1/4}=0{,}4172$.
\end{enumerate}
\resultado \textbf{$P(0{,}3)\approx 0{,}184$; $P(0{,}85)\approx 1-10^{-8}$; umbral aproximado $0{,}447$; inflexión $0{,}417$.}
\comprobacion $0{,}3<0{,}417<0{,}85$, así que la curva debe dar una probabilidad baja en el primer caso y casi~$1$ en el segundo. Con $r=4$, pequeño, la aproximación $(1/b)^{1/r}$ se separa algo de la inflexión exacta, como prevé la \cref{prop:curva-s}.
\end{solucion}

\begin{solucion}{eje:2-curvas}
\datos $(r,b)=(3,10)$, $(6,20)$, $(5,50)$; $s=0{,}2$, $0{,}5$, $0{,}8$.
\planteamiento $P(s)=1-(1-s^{r})^{b}$; de $P=1/2$ se despeja $s_{1/2}=(1-2^{-1/b})^{1/r}$.
\pasos
\begin{enumerate}[label=\arabic*.]
  \item $(3,10)$. $s=0{,}2$: $s^{3}=0{,}008$; $0{,}992^{10}=0{,}9228$; $P=0{,}0772$. $s=0{,}5$: $0{,}875^{10}=0{,}2631$; $P=0{,}7369$. $s=0{,}8$: $0{,}488^{10}=7{,}66\cdot 10^{-4}$; $P=0{,}9992$.
  \item $(6,20)$. $s=0{,}2$: $s^{6}=6{,}4\cdot 10^{-5}$; $P=0{,}0013$. $s=0{,}5$: $s^{6}=0{,}015625$; $0{,}984375^{20}=0{,}7298$; $P=0{,}2702$. $s=0{,}8$: $s^{6}=0{,}2621$; $0{,}7379^{20}=0{,}00229$; $P=0{,}9977$.
  \item $(5,50)$. $s=0{,}2$: $s^{5}=3{,}2\cdot 10^{-4}$; $0{,}99968^{50}=0{,}9841$; $P=0{,}0159$. $s=0{,}5$: $0{,}96875^{50}=0{,}2044$; $P=0{,}7956$. $s=0{,}8$: $0{,}6723^{50}=2{,}4\cdot 10^{-9}$; $P\approx 1$.
  \item $s_{1/2}$ para $(3,10)$: $2^{-1/10}=0{,}9330$; $0{,}0670^{1/3}=0{,}4061$, frente a $(1/10)^{1/3}=0{,}4642$.
  \item Para $(6,20)$: $2^{-1/20}=0{,}9659$; $0{,}0341^{1/6}=0{,}5694$, frente a $(1/20)^{1/6}=0{,}6070$.
  \item Para $(5,50)$: $2^{-1/50}=0{,}9862$; $0{,}0138^{1/5}=0{,}4244$, frente a $(1/50)^{1/5}=0{,}4573$.
\end{enumerate}
\resultado \textbf{(a) $(0{,}077;\ 0{,}737;\ 0{,}999)$ para $(3,10)$; $(0{,}001;\ 0{,}270;\ 0{,}998)$ para $(6,20)$; $(0{,}016;\ 0{,}796;\ 1{,}000)$ para $(5,50)$. (b) $s_{1/2}=0{,}406$; $0{,}569$; $0{,}424$, frente a $0{,}464$; $0{,}607$; $0{,}457$: la aproximación sobrestima el punto del $50\,\%$ entre $0{,}03$ y $0{,}06$.}
\comprobacion $0{,}4061^{3}=0{,}0670$ y $0{,}9330^{10}=0{,}500$. Además, $(6,20)$, con más filas por banda, es la curva más exigente: da la menor $P$ en $s=0{,}5$.
\end{solucion}

\begin{solucion}{eje:2-varianza}
\datos $\Var(\hat J)=J(1-J)/n$ (\cref{cor:estimador}).
\planteamiento acotar $J(1-J)$ en $[0,1]$; despejar $n$; aplicar Chebyshev, $\PP(\abs{X-\E X}\ge\varepsilon)\le\Var(X)/\varepsilon^{2}$ (módulo~5).
\pasos
\begin{enumerate}[label=\arabic*.]
  \item (a) $J(1-J)=\tfrac14-(J-\tfrac12)^{2}\le\tfrac14$, con igualdad en $J=1/2$. Luego $\sqrt{\Var(\hat J)}\le\sqrt{1/(4n)}=1/(2\sqrt n)$.
  \item (b) Para todo par: $1/(2\sqrt n)\le 0{,}05\iff\sqrt n\ge 10\iff n\ge 100$.
  \item Con $J=0{,}8$: $J(1-J)=0{,}16$; $0{,}16/n\le 0{,}0025\iff n\ge 64$.
  \item (c) $\PP(\abs{\hat J-J}\ge 0{,}1)\le\dfrac{J(1-J)}{100\cdot 0{,}01}\le\dfrac{0{,}25}{1}=0{,}25$.
\end{enumerate}
\resultado \textbf{(a) $\sigma\le 1/(2\sqrt n)$; (b) $n=100$ en general y $n=64$ si $J=0{,}8$; (c) probabilidad a lo sumo $0{,}25$.}
\comprobacion con $n=100$ y $J=1/2$, el error de $0{,}1$ son dos desviaciones típicas; la aproximación normal daría una probabilidad de unos $0{,}05$, compatible con la cota de Chebyshev, que es más conservadora.
\end{solucion}

% ---------------------------------------------------------------------
\section{Tema 3}

\begin{solucion}{eje:3-recta}
\datos $L=kN^{-\alpha}$; $\alpha=0{,}076$.
\planteamiento tomar logaritmos con $\log(xy)=\log x+\log y$ y $\log x^{a}=a\log x$ (\cref{prop:loglog}).
\pasos
\begin{enumerate}[label=\arabic*.]
  \item $\log L=\log k+\log N^{-\alpha}=\log k-\alpha\log N$.
  \item Con $u=\log N$ y $v=\log L$: $v=\log k-\alpha u$, una recta de pendiente $-\alpha$ y ordenada $\log k$.
  \item Multiplicar $N$ por $10$ aumenta $u=\log_{10}N$ en $1$, luego $v$ baja $\alpha=0{,}076$.
\end{enumerate}
\resultado \textbf{la ley es una recta de pendiente $-\alpha$; la pendiente es el cambio de $\log L$ por cada unidad de $\log N$. Con $\alpha=0{,}076$, $\log_{10}L$ baja $0{,}076$ por cada factor~$10$ en $N$.}
\comprobacion $10^{-0{,}076}=0{,}8395$: $L$ se multiplica por $0{,}84$, como en el \cref{ej:kaplan}.
\end{solucion}

\begin{solucion}{eje:3-computo}
\datos $N=1{,}3\cdot 10^{9}$; $D=2{,}6\cdot 10^{10}$; $1$ PF-día $=8{,}64\cdot 10^{19}$ FLOP.
\planteamiento $C\approx 6ND$; tokens por parámetro $=D/N$.
\pasos
\begin{enumerate}[label=\arabic*.]
  \item $6N=7{,}8\cdot 10^{9}$.
  \item $C=7{,}8\cdot 10^{9}\cdot 2{,}6\cdot 10^{10}=20{,}28\cdot 10^{19}=2{,}028\cdot 10^{20}$ FLOP.
  \item En PF-días: $2{,}028\cdot 10^{20}/8{,}64\cdot 10^{19}=2{,}347$.
  \item $D/N=2{,}6\cdot 10^{10}/1{,}3\cdot 10^{9}=20$.
\end{enumerate}
\resultado \textbf{$C=2{,}028\cdot 10^{20}$ FLOP $\approx 2{,}35$ PF-días; $20$ tokens por parámetro.}
\comprobacion con la \cref{prop:regla20}, $\sqrt{C/120}=\sqrt{1{,}69\cdot 10^{18}}=1{,}3\cdot 10^{9}$: se recupera el $N$ de partida.
\end{solucion}

\begin{solucion}{eje:3-datos}
\datos $L(D)=(5{,}4\cdot 10^{13}/D)^{0{,}095}$; $D=10^{9}$ y $10^{11}$.
\planteamiento $x^{0{,}095}=10^{0{,}095\log_{10}x}$.
\pasos
\begin{enumerate}[label=\arabic*.]
  \item $D=10^{9}$: $x=5{,}4\cdot 10^{4}$; $\log_{10}x=4{,}73239$; $0{,}095\cdot 4{,}73239=0{,}44958$; $L=10^{0{,}44958}=2{,}816$.
  \item $D=10^{11}$: $x=540$; $\log_{10}x=2{,}73239$; $0{,}095\cdot 2{,}73239=0{,}25958$; $L=10^{0{,}25958}=1{,}818$.
  \item Factor por cada factor~$10$: $10^{-0{,}095}=0{,}8035$.
\end{enumerate}
\resultado \textbf{$L(10^{9})\approx 2{,}816$ nats; $L(10^{11})\approx 1{,}818$ nats; cada factor~$10$ en $D$ multiplica $L$ por $0{,}804$.}
\comprobacion de $10^{9}$ a $10^{11}$ hay dos factores~$10$: $2{,}816\cdot 0{,}8035^{2}=2{,}816\cdot 0{,}6456=1{,}818$.
\end{solucion}

\begin{solucion}{eje:3-ajuste}
\datos configuraciones $(4,256)$, $(8,512)$, $(16,1024)$ con $L=4{,}20$, $3{,}62$, $3{,}12$; $N=12n_{\mathrm{layer}}d_{\mathrm{model}}^{2}$.
\planteamiento calcular $u_j=\log_{10}N_j$ y $v_j=\log_{10}L_j$ y aplicar la \cref{def:ajuste}.
\pasos
\begin{enumerate}[label=\arabic*.]
  \item (a) $N_1=12\cdot 4\cdot 65\,536=3\,145\,728$; $N_2=12\cdot 8\cdot 262\,144=25\,165\,824$; $N_3=12\cdot 16\cdot 1\,048\,576=201\,326\,592$. Cada uno es $8$ veces el anterior.
  \item $u=(6{,}49772;\ 7{,}40081;\ 8{,}30390)$, equiespaciados con paso $\log_{10}8=0{,}90309$; $\bar u=7{,}40081$.
  \item $v=(0{,}62325;\ 0{,}55871;\ 0{,}49415)$; $\bar v=0{,}55870$.
  \item Numerador: $(-0{,}90309)(0{,}06455)+0+(0{,}90309)(-0{,}06455)=-0{,}11658$. Denominador: $2\cdot 0{,}90309^{2}=1{,}63114$.
  \item $b_1=-0{,}11658/1{,}63114=-0{,}07147$; $b_0=0{,}55870+0{,}07147\cdot 7{,}40081=1{,}08767$.
  \item (c) $N=12\cdot 32\cdot 2048^{2}=1\,610\,612\,736$; $\log_{10}N=9{,}20700$; $\log_{10}\hat L=1{,}08767-0{,}07147\cdot 9{,}20700=0{,}42963$; $\hat L=2{,}689$.
\end{enumerate}
\resultado \textbf{(a) $N=3{,}15\cdot 10^{6}$, $2{,}52\cdot 10^{7}$, $2{,}01\cdot 10^{8}$; (b) $\hat L=10^{1{,}0877}N^{-0{,}0715}$, exponente $\hat\alpha=0{,}0715$; (c) $\hat L\approx 2{,}69$ nats para $N=1{,}61\cdot 10^{9}$.} La predicción extrapola un orden de magnitud con solo tres puntos y supone que los hiperparámetros (tasa de aprendizaje, lote, tokens) se ajustan bien en todos los tamaños; sin varias semillas no hay barras de error.
\comprobacion la recta ajustada reproduce los datos: $10^{1{,}08767-0{,}07147\cdot 6{,}49772}=4{,}200$, $3{,}620$ y $3{,}120$. El exponente es del orden del $\alpha_N=0{,}076$ de Kaplan \emph{et al.}
\end{solucion}

\begin{solucion}{eje:3-sobreajuste}
\datos $L(N,D)$ de~\eqref{eq:kaplan-nd}; $L(N)=(N_c/N)^{\alpha_N}$; $\alpha_N=0{,}076$; $\alpha_D=0{,}095$ o $0{,}103$.
\planteamiento sacar factor común $(N_c/N)^{\alpha_N/\alpha_D}$ dentro del corchete.
\pasos
\begin{enumerate}[label=\arabic*.]
  \item Sea $\rho=\alpha_N/\alpha_D$. El corchete es $(N_c/N)^{\rho}\bigl(1+(N/N_c)^{\rho}D_c/D\bigr)$.
  \item Elevando a $\alpha_D$ y usando $\rho\alpha_D=\alpha_N$:
  \[
  L(N,D)=\Bigl(\frac{N_c}{N}\Bigr)^{\rho\alpha_D}\Bigl(1+\Bigl(\frac{N}{N_c}\Bigr)^{\rho}\frac{D_c}{D}\Bigr)^{\alpha_D}=L(N)\Bigl(1+\frac{N^{\rho}D_c}{N_c^{\rho}D}\Bigr)^{\alpha_D},
  \]
  que es la igualdad~(a).
  \item (b) El cociente depende de $N$ y $D$ solo a través de $z=N^{\rho}/D$, y es creciente en $z$. Para mantenerlo constante al crecer $N$, $z$ debe mantenerse constante: $D\propto N^{\rho}=N^{\alpha_N/\alpha_D}$.
  \item (c) Con $\alpha_D=0{,}095$: $\rho=0{,}076/0{,}095=0{,}80$. Con $\alpha_D=0{,}103$: $\rho=0{,}076/0{,}103=0{,}738$.
\end{enumerate}
\resultado \textbf{$L(N,D)/L(N)=\bigl(1+(N/N_c)^{\alpha_N/\alpha_D}D_c/D\bigr)^{\alpha_D}$; $D$ debe crecer como $N^{\alpha_N/\alpha_D}$, con exponente $0{,}80$ o $0{,}74$ según el ajuste.} El segundo valor es el de la regla $D\gtrsim 5\cdot 10^{3}N^{0{,}74}$.
\comprobacion cuando $D\to\infty$, el cociente tiende a $1$, de acuerdo con la \cref{prop:limites}; el crecimiento sublineal ($\rho<1$) significa que los modelos grandes necesitan menos datos por parámetro para no sobreajustar.
\end{solucion}

\begin{solucion}{eje:3-invariancia}
\datos $b_1=\sum(u_j-\bar u)(v_j-\bar v)/\sum(u_j-\bar u)^{2}$, $b_0=\bar v-b_1\bar u$, $\hat\alpha=-b_1$.
\planteamiento cada transformación cambia $u$ o $v$ de forma afín; ver cómo cambian las desviaciones respecto de la media.
\pasos
\begin{enumerate}[label=\arabic*.]
  \item (a) Cambiar de base multiplica todos los logaritmos por la misma constante $\kappa_0>0$: $u'_j=\kappa_0u_j$, $v'_j=\kappa_0v_j$. Las desviaciones se multiplican por $\kappa_0$; numerador y denominador de $b_1$ por $\kappa_0^{2}$, y $b_1'=b_1$. La ordenada pasa a $b_0'=\kappa_0b_0$.
  \item (b) $N'=\kappa N$ da $u'_j=u_j+\log\kappa$: $\bar u'=\bar u+\log\kappa$ y $u'_j-\bar u'=u_j-\bar u$. Luego $b_1'=b_1$ y $b_0'=\bar v-b_1(\bar u+\log\kappa)=b_0-b_1\log\kappa$.
  \item (c) $L'=cL$ da $v'_j=v_j+\log c$, con $v'_j-\bar v'=v_j-\bar v$: $b_1'=b_1$ y $b_0'=b_0+\log c$. De nats a bits, $c=1/\ln 2$.
\end{enumerate}
\resultado \textbf{$\hat\alpha$ es invariante en los tres casos; cambia la ordenada: $b_0'=\kappa_0b_0$, $b_0'=b_0-b_1\log\kappa$ y $b_0'=b_0+\log c$, respectivamente.}
\comprobacion es coherente con la \cref{prop:escala}: el exponente es la única característica de una ley de potencia que no depende de las unidades, y la constante absorbe todos los cambios de escala.
\end{solucion}

% ---------------------------------------------------------------------
\section{Tema 4}

\begin{solucion}{eje:4-vf}
\planteamiento recordar los exponentes de $\Nopt$ en cada trabajo y el papel de $E$.
\pasos
\begin{enumerate}[label=(\alph*)]
  \item \textbf{Verdadera.} Kaplan \emph{et al.}\ obtienen $\Nopt\propto C^{0{,}73}$ y $D\propto C^{0{,}27}$: el exponente del modelo es casi el triple que el de los datos.
  \item \textbf{Falsa.} A cómputo fijo hay un tamaño óptimo (\cref{teo:chinchilla}); Gopher era demasiado grande para su presupuesto y Chinchilla, cuatro veces menor, obtiene menos pérdida: $1{,}937$ frente a $1{,}993$ (\cref{ej:gopher-chinchilla}).
  \item \textbf{Falsa.} $E=1{,}69$ no depende de $N$ ni de $D$: es la entropía del texto, y $L\to E$ cuando $N,D\to\infty$.
\end{enumerate}
\resultado \textbf{(a) V; (b) F; (c) F.}
\comprobacion en (c), con $N=D=10^{30}$ los términos de error valen $2{,}6\cdot 10^{-8}$ y $1{,}6\cdot 10^{-6}$, y $L\approx 1{,}69$.
\end{solucion}

\begin{solucion}{eje:4-regla20}
\datos $C=10^{23}$; $D=20N$.
\planteamiento por la \cref{prop:regla20}, $N=\sqrt{C/120}$.
\pasos
\begin{enumerate}[label=\arabic*.]
  \item $N^{2}=10^{23}/120=8{,}333\cdot 10^{20}$.
  \item $N=\sqrt{8{,}333}\cdot 10^{10}=2{,}887\cdot 10^{10}$.
  \item $D=20\cdot 2{,}887\cdot 10^{10}=5{,}774\cdot 10^{11}$.
\end{enumerate}
\resultado \textbf{$N\approx 2{,}89\cdot 10^{10}$ parámetros y $D\approx 5{,}77\cdot 10^{11}$ tokens.}
\comprobacion $6\cdot 2{,}887\cdot 10^{10}\cdot 5{,}774\cdot 10^{11}=1{,}000\cdot 10^{23}$.
\end{solucion}

\begin{solucion}{eje:4-optimo}
\datos $E=1{,}69$, $A=406{,}4$, $\alpha=0{,}34$, $B=410{,}7$, $\beta=0{,}28$; $C=10^{22}$.
\planteamiento \cref{teo:chinchilla}: $\Nopt=G(C/6)^{a}$, $\Dopt=(C/6)/\Nopt$; después $L=E+A\Nopt^{-\alpha}+B\Dopt^{-\beta}$.
\pasos
\begin{enumerate}[label=\arabic*.]
  \item $a=0{,}28/0{,}62=0{,}45161$; $1/(\alpha+\beta)=1{,}6129$.
  \item $\alpha A=138{,}18$; $\beta B=115{,}00$; cociente $1{,}20157$; $\ln 1{,}20157=0{,}18363$; $G=e^{1{,}6129\cdot 0{,}18363}=e^{0{,}29618}=1{,}3447$.
  \item $C/6=1{,}6667\cdot 10^{21}$; $\log_{10}(C/6)=21{,}22185$; $(C/6)^{a}=10^{0{,}45161\cdot 21{,}22185}=10^{9{,}58406}=3{,}8376\cdot 10^{9}$.
  \item $\Nopt=1{,}3447\cdot 3{,}8376\cdot 10^{9}=5{,}160\cdot 10^{9}$.
  \item $\Dopt=1{,}6667\cdot 10^{21}/5{,}160\cdot 10^{9}=3{,}230\cdot 10^{11}$; $\Dopt/\Nopt=62{,}6$.
  \item $\log_{10}\Nopt=9{,}71265$; $\Nopt^{0{,}34}=10^{3{,}30230}=2006$; $T_N=406{,}4/2006=0{,}2026$.
  \item $\log_{10}\Dopt=11{,}50920$; $\Dopt^{0{,}28}=10^{3{,}22258}=1669{,}5$; $T_D=410{,}7/1669{,}5=0{,}2460$.
  \item $L=1{,}69+0{,}2026+0{,}2460=2{,}1386$.
  \item Regla de~20: $N=\sqrt{10^{22}/120}=9{,}129\cdot 10^{9}$ y $D=1{,}826\cdot 10^{11}$. Términos de error: $406{,}4/10^{0{,}34\cdot 9{,}96042}=406{,}4/2435=0{,}1669$ y $410{,}7/10^{0{,}28\cdot 11{,}26150}=410{,}7/1423=0{,}2886$. Pérdida: $L=2{,}1455$.
\end{enumerate}
\resultado \textbf{$\Nopt\approx 5{,}16\cdot 10^{9}$; $\Dopt\approx 3{,}23\cdot 10^{11}$; $\Dopt/\Nopt\approx 63$; $L\approx 2{,}1386$ nats. La regla de~20 da $N\approx 9{,}13\cdot 10^{9}$, $D\approx 1{,}83\cdot 10^{11}$ y $L\approx 2{,}1455$: solo $0{,}007$ nats peor.}
\comprobacion (i) $6\cdot 5{,}160\cdot 10^{9}\cdot 3{,}230\cdot 10^{11}=1{,}000\cdot 10^{22}$. (ii) Por el \cref{cor:equilibrio}, $T_D/T_N$ debe valer $\alpha/\beta=1{,}214$: $0{,}2460/0{,}2026=1{,}214$.
\end{solucion}

\begin{solucion}{eje:4-plano}
\datos $T_N=0{,}2026$, $T_D=0{,}2460$, $\alpha=0{,}34$, $\beta=0{,}28$; factores $\lambda=4$ y $\lambda=1/4$.
\planteamiento \cref{prop:plano} con $\delta=\ln\lambda=\pm 1{,}3863$.
\pasos
\begin{enumerate}[label=\arabic*.]
  \item Coeficiente cuadrático: $\tfrac12(\alpha^{2}T_N+\beta^{2}T_D)=\tfrac12(0{,}1156\cdot 0{,}2026+0{,}0784\cdot 0{,}2460)=0{,}02135$; con $\delta^{2}=1{,}9218$, aproximación $0{,}0410$ en ambos casos.
  \item $\lambda=4$: $e^{-0{,}34\cdot 1{,}3863}=e^{-0{,}4713}=0{,}6242$ y $e^{0{,}28\cdot 1{,}3863}=e^{0{,}3882}=1{,}4743$. $\Delta=0{,}2026\cdot(-0{,}3758)+0{,}2460\cdot 0{,}4743=-0{,}07614+0{,}11668=0{,}0405$.
  \item $\lambda=1/4$: $e^{0{,}4713}=1{,}6021$ y $e^{-0{,}3882}=0{,}6783$. $\Delta=0{,}2026\cdot 0{,}6021+0{,}2460\cdot(-0{,}3217)=0{,}12198-0{,}07914=0{,}0429$.
  \item Asimetría: el término cúbico del desarrollo es $\tfrac16(-\alpha^{3}T_N+\beta^{3}T_D)\delta^{3}$, con $-\alpha^{3}T_N+\beta^{3}T_D=-0{,}00796+0{,}00540=-0{,}00256<0$; resta en $\delta>0$ y suma en $\delta<0$.
\end{enumerate}
\resultado \textbf{exceso exacto $0{,}0405$ nats con $4\Nopt$ y $0{,}0429$ con $\Nopt/4$; la aproximación cuadrática da $0{,}0410$ en ambos.} Equivocarse por exceso en el tamaño cuesta algo menos que por defecto, porque $\alpha>\beta$.
\comprobacion $\delta^{3}/6\cdot(-0{,}00256)=\mp 0{,}00114$, y $0{,}0410-0{,}0011=0{,}0399$ y $0{,}0410+0{,}0011=0{,}0421$ se acercan a los valores exactos; el resto lo aportan los términos de orden~4.
\end{solucion}

\begin{solucion}{eje:4-factor4}
\datos $\Nopt=G(C/6)^{a}$, $a=\beta/(\alpha+\beta)$; $\Dopt=(C/6)/\Nopt$; $\alpha=0{,}34$, $\beta=0{,}28$.
\planteamiento dividir la fórmula en $4C$ entre la fórmula en $C$; $G$ y el factor $1/6$ se cancelan.
\pasos
\begin{enumerate}[label=\arabic*.]
  \item $\Nopt(4C)/\Nopt(C)=G(4C/6)^{a}/\bigl(G(C/6)^{a}\bigr)=4^{a}$.
  \item $\Dopt=G^{-1}(C/6)^{1-a}$, porque $x/x^{a}=x^{1-a}$; del mismo modo, $\Dopt(4C)/\Dopt(C)=4^{1-a}$.
  \item $1-a=(\alpha+\beta-\beta)/(\alpha+\beta)=\alpha/(\alpha+\beta)$.
  \item $a=0{,}45161$ y $1-a=0{,}54839$; $\log_{10}4=0{,}60206$.
  \item $4^{0{,}45161}=10^{0{,}27190}=1{,}870$; $4^{0{,}54839}=10^{0{,}33017}=2{,}139$.
\end{enumerate}
\resultado \textbf{$\Nopt$ se multiplica por $4^{\beta/(\alpha+\beta)}\approx 1{,}870$ y $\Dopt$ por $4^{\alpha/(\alpha+\beta)}\approx 2{,}139$.}
\comprobacion $1{,}870\cdot 2{,}139=4{,}000$, como exige $6\Nopt\Dopt=C$.
\end{solucion}

\begin{solucion}{eje:4-perdida}
\datos $\Nopt=Gc^{a}$, $\Dopt=G^{-1}c^{b}$, $c=C/6$; valores de Chinchilla; del \cref{eje:4-optimo}, $T_N+T_D=0{,}4486$ en $10^{22}$.
\planteamiento seguir la demostración del \cref{cor:perdida} conservando las constantes.
\pasos
\begin{enumerate}[label=\arabic*.]
  \item $A\Nopt^{-\alpha}=AG^{-\alpha}c^{-\gamma}$ y $B\Dopt^{-\beta}=BG^{\beta}c^{-\gamma}$, con $\gamma=\alpha a=\beta b=\alpha\beta/(\alpha+\beta)$.
  \item $L^{*}-E=(AG^{-\alpha}+BG^{\beta})c^{-\gamma}=(AG^{-\alpha}+BG^{\beta})\,6^{\gamma}\,C^{-\gamma}$, luego $K=(AG^{-\alpha}+BG^{\beta})\,6^{\gamma}$.
  \item $\gamma=0{,}0952/0{,}62=0{,}15355$. $G^{-\alpha}=1{,}3447^{-0{,}34}=0{,}9042$ y $G^{\beta}=1{,}3447^{0{,}28}=1{,}0865$; $AG^{-\alpha}=367{,}47$ y $BG^{\beta}=446{,}21$; suma $813{,}68$; $6^{\gamma}=1{,}3167$; $K=1071{,}4$.
  \item En $10^{22}$: $C^{-\gamma}=10^{-22\cdot 0{,}15355}=4{,}187\cdot 10^{-4}$, y $K C^{-\gamma}=0{,}4486$.
  \item De $10^{22}$ a $10^{24}$: factor $100^{-\gamma}=10^{-0{,}30710}=0{,}4931$; exceso $0{,}4486\cdot 0{,}4931=0{,}2212$.
  \item Directamente en $10^{24}$: $\Nopt=4{,}130\cdot 10^{10}$, $\Dopt=4{,}036\cdot 10^{12}$, $T_N=0{,}0999$, $T_D=0{,}1213$; suma $0{,}2212$.
\end{enumerate}
\resultado \textbf{$K=(AG^{-\alpha}+BG^{\beta})6^{\gamma}\approx 1071$, $\gamma\approx 0{,}154$. Con $100$ veces más cómputo, el exceso pasa de $0{,}449$ a $0{,}221$ nats (factor $0{,}493$).}
\comprobacion el cálculo directo del paso~6 coincide con la ley del paso~5, y el valor de $K C^{-\gamma}$ en $10^{22}$ reproduce la suma $T_N+T_D$ del \cref{eje:4-optimo}.
\end{solucion}

\begin{solucion}{eje:4-gp}
\datos $f(N,D)=AN^{-\alpha}+BD^{-\beta}$; restricción $6ND\le C$; $T_N=0{,}2026$, $T_D=0{,}2460$, $\gamma=0{,}15355$ en $C=10^{22}$.
\planteamiento cambio de variables logarítmico~\cite[§4.5]{bv}; condiciones KKT~\cite[§5.5]{bv}; comparación con el \cref{cor:perdida}.
\pasos
\begin{enumerate}[label=\arabic*.]
  \item (a) Con $N=e^{u}$, $D=e^{v}$: minimizar $\tilde f(u,v)=Ae^{-\alpha u}+Be^{-\beta v}$ sujeto a $u+v-\ln(C/6)\le 0$. Cada sumando es la exponencial, convexa y creciente, de una función afín, luego convexo; la suma es convexa. La restricción es afín. El problema es convexo.
  \item (b) Si en un punto factible la restricción no fuera activa, $u+v<\ln(C/6)$, aumentar $u$ ligeramente seguiría siendo factible y reduciría estrictamente $Ae^{-\alpha u}$. Luego en el óptimo $u+v=\ln(C/6)$, es decir, $6ND=C$.
  \item (c) En las variables originales, KKT: $\nabla f+\mu\nabla(6ND-C)=0$, $\mu\ge 0$, $\mu(6ND-C)=0$. Las componentes dan $-\alpha AN^{-\alpha-1}+6\mu D=0$ y $-\beta BD^{-\beta-1}+6\mu N=0$, que son las del \cref{teo:chinchilla} con $\lambda=\mu$. De la primera, $\mu=\alpha AN^{-\alpha}/(6ND)=\alpha T_N/C>0$, compatible con la restricción activa. Como el problema transformado es convexo, estas condiciones son suficientes.
  \item Numéricamente, $\mu C=\alpha T_N=0{,}34\cdot 0{,}2026=0{,}0689$.
  \item (d) Por el \cref{cor:perdida}, $dL^{*}/d\ln C=-\gamma KC^{-\gamma}=-\gamma(T_N+T_D)$. Como $T_D=(\alpha/\beta)T_N$, $\gamma(T_N+T_D)=\gamma T_N(\alpha+\beta)/\beta=\alpha T_N=\mu C$. Luego $dL^{*}/d\ln C=-\mu C$, es decir, $dL^{*}/dC=-\mu$.
  \item Con los datos: $\gamma(T_N+T_D)=0{,}15355\cdot 0{,}4486=0{,}0689$.
\end{enumerate}
\resultado \textbf{el problema es convexo en $(\ln N,\ln D)$; la restricción es activa; el punto del \cref{teo:chinchilla} satisface KKT con $\mu=\alpha T_N/C>0$, y $\mu$ es el ritmo al que baja la pérdida óptima por unidad de cómputo: $dL^{*}/d\ln C=-0{,}0689$ nats en $10^{22}$ FLOP.}
\comprobacion duplicar el cómputo ($\Delta\ln C=0{,}693$) debería reducir la pérdida en unos $0{,}0689\cdot 0{,}693=0{,}048$ nats; el \cref{cor:perdida} da $0{,}4486\cdot(1-2^{-0{,}15355})=0{,}4486\cdot 0{,}1009=0{,}0453$, del mismo orden (la diferencia procede de la curvatura).
\end{solucion}

% ---------------------------------------------------------------------
\section{Tema 5}

\begin{solucion}{eje:5-sobreentreno}
\datos coste de entrenamiento $6ND$; coste de inferencia $2N$ por token servido.
\planteamiento comparar costes totales (\cref{prop:amortizacion}).
\pasos
\begin{enumerate}[label=\arabic*.]
  \item El óptimo de Chinchilla minimiza solo $6ND$ para una pérdida dada. Un modelo menor necesita más tokens para igualar la pérdida, y su entrenamiento cuesta más (\cref{prop:sobrecoste}).
  \item Pero cada token servido cuesta $2N$: un modelo la mitad de grande ahorra la mitad del cómputo de inferencia, para siempre. Con millones de usuarios, $T_{\mathrm{inf}}$ supera enseguida el umbral $T^{*}$ de la \cref{prop:amortizacion}.
  \item Además, un modelo menor requiere menos memoria y responde con menos latencia.
  \item El óptimo de Chinchilla es preferible cuando el modelo se usará poco (un experimento, un prototipo, un modelo de investigación) o cuando el cuello de botella es el cómputo de entrenamiento disponible.
\end{enumerate}
\resultado \textbf{porque el coste total incluye $2N T_{\mathrm{inf}}$, que domina con mucha demanda; el óptimo de Chinchilla conviene cuando se servirán pocos tokens.}
\comprobacion en el \cref{ej:sobreentreno}, el sobrecoste se amortiza tras $1{,}1\cdot 10^{11}$ tokens: $10^{6}$ usuarios que generan $10^{5}$ tokens cada uno ya lo superan.
\end{solucion}

\begin{solucion}{eje:5-isoflop}
\datos la tabla de nueve entrenamientos en tres presupuestos.
\planteamiento mínimo de cada perfil; recta en $(\log_{10}C,\log_{10}\Nopt)$; $D=C/(6N)$.
\pasos
\begin{enumerate}[label=\arabic*.]
  \item Mínimos: $10^{18}\to 5\cdot 10^{7}$ ($3{,}85$); $10^{19}\to 1{,}6\cdot 10^{8}$ ($3{,}40$); $10^{20}\to 5\cdot 10^{8}$ ($3{,}01$).
  \item $x=18,19,20$; $y=7{,}69897;\ 8{,}20412;\ 8{,}69897$; $\bar x=19$; $\bar y=24{,}60206/3=8{,}20069$.
  \item Pendiente: numerador $(-1)(-0{,}50172)+0+(1)(0{,}49828)=1{,}00000$; denominador $2$; $a=0{,}5$.
  \item Ordenada: $b_0=8{,}20069-0{,}5\cdot 19=-1{,}29931$.
  \item $C=10^{23}$: $\log_{10}N=10{,}20069$; $N=1{,}587\cdot 10^{10}$; $D=10^{23}/(9{,}522\cdot 10^{10})=1{,}050\cdot 10^{12}$; $D/N=66{,}1$.
  \item $C=10^{24}$: $\log_{10}N=10{,}70069$; $N=5{,}020\cdot 10^{10}$; $D=3{,}320\cdot 10^{12}$; $D/N=66{,}1$.
\end{enumerate}
\resultado \textbf{$\Nopt\propto C^{0{,}5}$. En $10^{23}$: $N\approx 1{,}59\cdot 10^{10}$, $D\approx 1{,}05\cdot 10^{12}$; en $10^{24}$: $N\approx 5{,}02\cdot 10^{10}$, $D\approx 3{,}32\cdot 10^{12}$; unos $66$ tokens por parámetro, más que la regla de~20 en estos datos sintéticos.}
\comprobacion con $a=0{,}5$, $N$ y $D$ crecen como $\sqrt C$ y $D/N$ es constante: sale $66$ en ambos presupuestos; además, $6\cdot 1{,}587\cdot 10^{10}\cdot 1{,}050\cdot 10^{12}=1{,}000\cdot 10^{23}$.
\end{solucion}

\begin{solucion}{eje:5-vertice}
\datos $C=10^{19}$; $x=(8{,}0;\ 8{,}3;\ 8{,}6)$, $h=0{,}3$; $L=(3{,}47;\ 3{,}40;\ 3{,}43)$.
\planteamiento \cref{prop:vertice}; $D=C/(6N)$.
\pasos
\begin{enumerate}[label=\arabic*.]
  \item $\Delta_2=3{,}47-6{,}80+3{,}43=0{,}10>0$; $L_3-L_1=-0{,}04$.
  \item $x^{*}=8{,}3-0{,}3\cdot(-0{,}04)/(2\cdot 0{,}10)=8{,}3+0{,}06=8{,}36$; $\Nopt=10^{8{,}36}=2{,}291\cdot 10^{8}$.
  \item $L^{*}=3{,}40-0{,}0016/0{,}8=3{,}398$.
  \item $\Dopt=10^{19}/(6\cdot 2{,}291\cdot 10^{8})=7{,}275\cdot 10^{9}$; $\Dopt/\Nopt=31{,}8$.
  \item El criterio del mínimo observado daría $10^{8{,}3}=1{,}995\cdot 10^{8}$, un $13\,\%$ menos.
\end{enumerate}
\resultado \textbf{vértice en $\log_{10}N=8{,}36$: $\Nopt\approx 2{,}29\cdot 10^{8}$, $L^{*}\approx 3{,}398$, $\Dopt\approx 7{,}3\cdot 10^{9}$ ($32$ tokens por parámetro). El mínimo observado da $2{,}0\cdot 10^{8}$.}
\comprobacion la parábola $p(t)=3{,}40-0{,}02t+0{,}05t^{2}$ pasa por los tres puntos ($p(-1)=3{,}47$, $p(1)=3{,}43$) y su mínimo, en $t^{*}=0{,}2$, vale $3{,}40-0{,}004+0{,}002=3{,}398$. El vértice se desplaza hacia el lado de la pérdida menor ($L_3<L_1$).
\end{solucion}

\begin{solucion}{eje:5-repetir}
\datos $U=10^{11}$; $R^{*}=15{,}39$; entrenamientos de $4\cdot 10^{11}$ y $1{,}6\cdot 10^{12}$ tokens.
\planteamiento épocas $=D/U$; $R=$ épocas $-1$; $D'=U\bigl(1+R^{*}(1-e^{-R/R^{*}})\bigr)$.
\pasos
\begin{enumerate}[label=\arabic*.]
  \item Primer entrenamiento: $4$ épocas, $R=3$; $R/R^{*}=0{,}1949$; $1-e^{-0{,}1949}=0{,}1771$; $D'=10^{11}(1+15{,}39\cdot 0{,}1771)=3{,}726\cdot 10^{11}$; fracción $3{,}726/4=0{,}931$.
  \item Segundo: $16$ épocas, $R=15$; $R/R^{*}=0{,}9748$; $1-e^{-0{,}9748}=0{,}6227$; $D'=10^{11}(1+15{,}39\cdot 0{,}6227)=1{,}058\cdot 10^{12}$; fracción $1{,}058/1{,}6=0{,}661$.
  \item Máximo: $U(1+R^{*})=1{,}639\cdot 10^{12}$ (\cref{prop:efectivos}).
\end{enumerate}
\resultado \textbf{$D'=3{,}73\cdot 10^{11}$ ($93\,\%$ de lo procesado) con $4$ épocas y $D'=1{,}06\cdot 10^{12}$ ($66\,\%$) con $16$; el máximo alcanzable es $1{,}64\cdot 10^{12}$ tokens efectivos.}
\comprobacion cuadruplicar el entrenamiento de $4$ a $16$ épocas solo multiplica los datos efectivos por $1{,}058/0{,}3726=2{,}84$: el rendimiento marginal decrece, como exige la concavidad de la \cref{prop:efectivos}.
\end{solucion}

\begin{solucion}{eje:5-reto}
\datos $C=10^{22}$; del \cref{eje:4-optimo}, $N_1=5{,}1605\cdot 10^{9}$, $L_1=2{,}13861$; $N_2=N_1/2$; inferencia $2N$ por token.
\planteamiento despejar $D_2$ de $L(N_2,D_2)=L_1$; $C_2=6N_2D_2$; $T^{*}$ de la \cref{prop:amortizacion} con $6N_1D_1=C$. Se usan cinco cifras porque el exponente $1/0{,}28=3{,}571$ amplifica los errores.
\pasos
\begin{enumerate}[label=\arabic*.]
  \item $N_2=2{,}5802\cdot 10^{9}$; $\log_{10}N_2=9{,}41165$; $N_2^{0{,}34}=10^{3{,}19996}=1584{,}7$; $A/N_2^{\alpha}=0{,}25645$.
  \item $B/D_2^{\beta}=2{,}13861-1{,}69-0{,}25645=0{,}19216$.
  \item $D_2^{0{,}28}=410{,}7/0{,}19216=2137{,}3$; $\log_{10}D_2=3{,}32986/0{,}28=11{,}89236$; $D_2=7{,}803\cdot 10^{11}$.
  \item Tokens por parámetro: $7{,}803\cdot 10^{11}/2{,}5802\cdot 10^{9}=302$.
  \item $C_2=6\cdot 2{,}5802\cdot 10^{9}\cdot 7{,}803\cdot 10^{11}=1{,}208\cdot 10^{22}$; sobrecoste $2{,}08\cdot 10^{21}$ FLOP.
  \item Ahorro por token servido: $2(N_1-N_2)=5{,}160\cdot 10^{9}$ FLOP; $T^{*}=2{,}08\cdot 10^{21}/5{,}160\cdot 10^{9}=4{,}03\cdot 10^{11}$ tokens.
  \item \cref{prop:sobrecoste} con $\kappa=1/2$: $C_2/C=1{,}208$, igual que en el paso~5.
\end{enumerate}
\resultado \textbf{$D_2\approx 7{,}80\cdot 10^{11}$ tokens ($302$ por parámetro); el entrenamiento cuesta $1{,}208$ veces más ($+2{,}08\cdot 10^{21}$ FLOP) y se amortiza tras unos $4{,}0\cdot 10^{11}$ tokens de inferencia.}
\comprobacion sustituyendo, $1{,}69+0{,}25645+410{,}7/2137{,}3=1{,}69+0{,}25645+0{,}19216=2{,}13861=L_1$. El sobrecoste relativo es el mismo que a $10^{21}$ FLOP (\cref{ej:sobreentreno}), como predice la \cref{prop:sobrecoste}; el umbral de amortización, en cambio, escala con el presupuesto.
\end{solucion}

\begin{solucion}{eje:5-cota}
\datos $D'=U+UR^{*}(1-e^{-R/R^{*}})$, $R\ge 0$.
\planteamiento acotar $1-e^{-x}$ entre $x-x^{2}/2$ y $x$ para $x\ge 0$.
\pasos
\begin{enumerate}[label=\arabic*.]
  \item Sea $\psi(x)=x-(1-e^{-x})$. $\psi(0)=0$ y $\psi'(x)=1-e^{-x}\ge 0$: $1-e^{-x}\le x$.
  \item Sea $\chi(x)=(1-e^{-x})-x+x^{2}/2$. $\chi(0)=0$ y $\chi'(x)=e^{-x}-1+x\ge 0$ por el paso~1 aplicado a $e^{-x}\ge 1-x$: $1-e^{-x}\ge x-x^{2}/2$.
  \item Con $x=R/R^{*}$ y multiplicando por $UR^{*}>0$: $UR-UR^{2}/(2R^{*})\le UR^{*}(1-e^{-R/R^{*}})\le UR$. Sumando $U$ se obtiene la doble desigualdad.
  \item Dividiendo por $U(1+R)$: $1-D'/(U(1+R))\le R^{2}/\bigl(2R^{*}(1+R)\bigr)\le R/(2R^{*})$, porque $R/(1+R)\le 1$.
  \item Cuatro épocas ($R=3$): cota $3/(2\cdot 15{,}39)=0{,}0975$ (la intermedia, $9/(2\cdot 15{,}39\cdot 4)=0{,}0731$); valor exacto $1-0{,}9314=0{,}0686$.
\end{enumerate}
\resultado \textbf{$U(1+R)-UR^{2}/(2R^{*})\le D'\le U(1+R)$; la pérdida de valor es a lo sumo $R/(2R^{*})$: un $9{,}8\,\%$ con cuatro épocas, frente al $6{,}9\,\%$ exacto.}
\comprobacion con $R=3$, la cota inferior da $4-0{,}2924=3{,}708$ (en unidades de $U$), y el valor exacto, $3{,}726$, está entre ella y $4$.
\end{solucion}

% ---------------------------------------------------------------------
\section{Tema 6}

\begin{solucion}{eje:6-escepticismo}
\datos una afirmación de emergencia a partir de $10^{10}$ parámetros.
\planteamiento aplicar el argumento de la métrica (\cref{prop:gompertz}) y las precauciones experimentales habituales.
\pasos
\begin{enumerate}[label=\arabic*.]
  \item \emph{La métrica.} Si es exactitud exacta sobre respuestas de $k$ tokens, vale $p^{k}$: con $k=10$, pasar de $p=0{,}8$ a $0{,}95$ lleva la exactitud de $0{,}107$ a $0{,}599$. Hay que repetir la medida con una métrica continua (log-probabilidad de la respuesta correcta, distancia de edición). Si la curva es suave, el salto lo crea la métrica.
  \item \emph{La resolución.} Con tres o cuatro tamaños de modelo, cualquier curva en S parece un escalón. Hacen falta tamaños intermedios y varias semillas para estimar barras de error.
  \item \emph{La contaminación.} Un modelo mayor memoriza más; si la tarea aparece en el corpus, puede «aprobar» por haber visto las respuestas.
\end{enumerate}
\resultado \textbf{comprobaría la métrica (todo o nada frente a continua), la densidad de tamaños con su incertidumbre y la contaminación del corpus.}
\comprobacion con la tabla del \cref{tema:6}, pasar de $p=0{,}8$ a $0{,}9$ con $k=10$ multiplica la exactitud por $0{,}349/0{,}107\approx 3{,}3$, mientras que $p$ solo crece un $12{,}5\,\%$.
\end{solucion}

\begin{solucion}{eje:6-todo-nada}
\datos $k$ tokens independientes con probabilidad de acierto $p$.
\planteamiento exactitud $p^{k}$ (independencia); número esperado de aciertos $kp$ (esperanza de una binomial).
\pasos
\begin{enumerate}[label=\arabic*.]
  \item (a) $p^{10}=0{,}5\iff p=0{,}5^{1/10}=0{,}9330$. Con $k=20$: $p=0{,}5^{1/20}=0{,}9659$.
  \item (b) $0{,}9^{10}=0{,}3487$ y $0{,}8^{10}=0{,}1074$; cociente $(0{,}9/0{,}8)^{10}=1{,}125^{10}=3{,}247$.
  \item (c) $kp$ pasa de $8$ a $9$: un $12{,}5\,\%$ más, el mismo aumento relativo que $p$.
\end{enumerate}
\resultado \textbf{(a) $p\approx 0{,}933$ con $k=10$ y $p\approx 0{,}966$ con $k=20$; (b) la exactitud se multiplica por $3{,}25$; (c) los aciertos esperados crecen un $12{,}5\,\%$ (de $8$ a $9$).}
\comprobacion $0{,}9330^{10}=0{,}500$; la respuesta más larga exige una $p$ mayor para la misma exactitud, coherente con el desplazamiento del umbral de la \cref{prop:gompertz}.
\end{solucion}

\begin{solucion}{eje:6-atencion}
\datos $N=1{,}3\cdot 10^{10}$; $L=40$; $d=5120$.
\planteamiento coste por token $2N+4L\nctx d$ (módulo~14); umbral $n^{*}=N/(2Ld)$.
\pasos
\begin{enumerate}[label=\arabic*.]
  \item (a) $n^{*}=1{,}3\cdot 10^{10}/(2\cdot 40\cdot 5120)=1{,}3\cdot 10^{10}/409\,600=31\,738$ tokens.
  \item (b) $\nctx=8192$: $4\cdot 40\cdot 8192\cdot 5120=6{,}711\cdot 10^{9}$; $2N=2{,}6\cdot 10^{10}$; fracción $6{,}711/(26+6{,}711)=0{,}205$.
  \item $\nctx=131\,072$: $1{,}074\cdot 10^{11}$; fracción $1{,}074/(0{,}26+1{,}074)=0{,}805$.
\end{enumerate}
\resultado \textbf{(a) a partir de unos $31\,700$ tokens; (b) la atención es el $20{,}5\,\%$ del coste por token con $8192$ tokens y el $80{,}5\,\%$ con $131\,072$.}
\comprobacion el término de atención es lineal en $\nctx$: $131\,072/8192=16$ y $16\cdot 6{,}711\cdot 10^{9}=1{,}074\cdot 10^{11}$.
\end{solucion}

\begin{solucion}{eje:6-kv}
\datos $L=40$; $h_{kv}=40$, $d_k=128$ ($h_{kv}d_k=5120$); bf16, $\beta_b=2$.
\planteamiento \cref{prop:kv}; $1$ MiB $=2^{20}$ bytes y $1$ GiB $=2^{30}$ bytes.
\pasos
\begin{enumerate}[label=\arabic*.]
  \item (a) Por token: $2\cdot 40\cdot 5120\cdot 2=819\,200$ bytes $=0{,}78125$ MiB.
  \item (b) $32\,768$ tokens: $819\,200\cdot 32\,768=26\,843\,545\,600$ bytes $=25$ GiB. Lote de $8$: $200$ GiB.
  \item (c) Con $h_{kv}=8$, la caché se multiplica por $8/40=1/5$: $5$ GiB por secuencia.
\end{enumerate}
\resultado \textbf{(a) $0{,}78$ MiB por token; (b) $25$ GiB por secuencia y $200$ GiB por lote; (c) $5$ GiB con GQA de $8$ cabezas.}
\comprobacion $819\,200=0{,}78125\cdot 2^{20}$ y $0{,}78125\cdot 32\,768=25\,600$ MiB $=25$ GiB: la caché de una sola secuencia larga supera ya la mitad de una GPU de $40$ GiB.
\end{solucion}

\begin{solucion}{eje:6-gompertz}
\datos $\ell(N)=(10^{8}/N)^{1/2}$; $A(N)=e^{-k\ell(N)}$.
\planteamiento evaluar $A$; usar la \cref{prop:gompertz}~(ii) para la inflexión; derivar $A$ respecto de $x=\ln N$.
\pasos
\begin{enumerate}[label=\arabic*.]
  \item (a) $k=5$: $N=10^{8}$: $\ell=1$, $A=e^{-5}=0{,}0067$. $N=10^{9}$: $\ell=0{,}3162$, $A=e^{-1{,}581}=0{,}2057$. $N=10^{10}$: $\ell=0{,}1$, $A=e^{-0{,}5}=0{,}6065$. $N=10^{11}$: $\ell=0{,}0316$, $A=e^{-0{,}1581}=0{,}8538$.
  \item (b) $N_0=N_1k^{1/\theta}=10^{8}k^{2}$: $10^{8}$, $2{,}5\cdot 10^{9}$ y $6{,}25\cdot 10^{10}$ para $k=1$, $5$ y $25$ ($\log_{10}N_0=8$; $9{,}40$; $10{,}80$).
  \item Con $k\to\lambda k$: $\log_{10}N_0$ pasa de $\log_{10}N_1+(\log_{10}k)/\theta$ a $\log_{10}N_1+(\log_{10}k+\log_{10}\lambda)/\theta$: se desplaza $(\log_{10}\lambda)/\theta$ órdenes de magnitud. Con $\lambda=5$ y $\theta=1/2$, $2\log_{10}5=1{,}40$.
  \item (c) Por la demostración de la \cref{prop:gompertz}, $dA/dx=\theta ye^{-y}$ con $y=e^{-\theta(x-x_0)}>0$. La función $ye^{-y}$ alcanza su máximo, $e^{-1}$, en $y=1$; luego $\max dA/dx=\theta/e$, que no depende de $k$.
\end{enumerate}
\resultado \textbf{(a) $A=0{,}007$; $0{,}206$; $0{,}607$; $0{,}854$; (b) inflexiones en $10^{8}$, $2{,}5\cdot 10^{9}$ y $6{,}25\cdot 10^{10}$; el umbral se desplaza $(\log_{10}\lambda)/\theta$ órdenes de magnitud; (c) pendiente máxima $\theta/e$.} La «brusquedad» de la transición en escala logarítmica es la misma para cualquier longitud de respuesta: alargar la respuesta solo traslada la curva hacia modelos mayores.
\comprobacion en $N_0=2{,}5\cdot 10^{9}$ con $k=5$: $\ell=(1/25)^{1/2}=0{,}2$ y $A=e^{-1}=0{,}368$, como exige la \cref{prop:gompertz}; el valor cae entre los de $10^{9}$ y $10^{10}$ del apartado~(a).
\end{solucion}

\begin{solucion}{eje:6-pi}
\datos base $10^{4}$; $d_k=128$; $L_{\mathrm{tr}}=4096$; $L'=32\,768$.
\planteamiento frecuencia más baja $\theta_{63}=\text{base}^{-126/128}$; \cref{prop:pi}; para la base nueva, igualar $L'\theta'_{63}=L_{\mathrm{tr}}\theta_{63}$.
\pasos
\begin{enumerate}[label=\arabic*.]
  \item (a) $\theta_{63}=10^{-4\cdot 0{,}984375}=10^{-3{,}9375}=1{,}1548\cdot 10^{-4}$. Ángulo máximo en el entrenamiento: $4096\cdot 1{,}1548\cdot 10^{-4}=0{,}4730$ rad. Sin interpolar, en $L'$: $32\,768\cdot 1{,}1548\cdot 10^{-4}=3{,}784$ rad.
  \item (b) Con interpolación, la posición $m<L'$ se sustituye por $m/8$, y el ángulo es $\theta_i m/8<\theta_iL'/8=\theta_iL_{\mathrm{tr}}$, dentro del rango visto (\cref{prop:pi}).
  \item (c) Se busca $b'$ con $L'(b')^{-126/128}=L_{\mathrm{tr}}\,10^{-4\cdot 126/128}$, es decir, $(b'/10^{4})^{126/128}=L'/L_{\mathrm{tr}}=8$, de donde $b'=10^{4}\cdot 8^{128/126}=10^{4}\cdot 8{,}2685=82\,685$.
  \item Para $i\ge 1$, $\theta_i'=(b')^{-2i/d_k}<10^{-4\cdot 2i/d_k}=\theta_i$ porque $b'>10^{4}$ y el exponente $-2i/d_k$ es negativo: todas las frecuencias giran más despacio.
\end{enumerate}
\resultado \textbf{(a) $0{,}473$ rad en el entrenamiento frente a $3{,}78$ rad sin interpolar; (b) con interpolación ningún ángulo supera el rango visto; (c) $b'\approx 82\,700$, y con ella todas las frecuencias giran más despacio que con la base original.}
\comprobacion $32\,768\cdot 82\,685^{-126/128}=32\,768\cdot 1{,}4436\cdot 10^{-5}=0{,}4730$ rad, el máximo del entrenamiento. A diferencia de la interpolación, aumentar la base no comprime por igual todas las frecuencias: las altas, que ya daban muchas vueltas en el entrenamiento, cambian poco.
\end{solucion}
